% ======================================================================
\section{Introduction}\label{sec:intro}

Over the past three decades, the world made unprecedented progress on basic human needs. In the countries studied in this paper, mean life expectancy rose by \nLeGain{} years between \nShareYearStart{} and \nLastYear, and median under-5 mortality fell from \nDescUFirst{} to \nDescULast{} deaths per 1{,}000 live births between 1990--2000 and 2011--2022. Over the same period, the pressures that human activity places on the Earth system intensified, and annual global \COtwo{} emissions rose well above their 1990 level \citep{friedlingstein2026}. By the 2023 assessment, six of the nine planetary boundaries that define a ``safe operating space'' for humanity had been crossed \citep{rockstrom2009,steffen2015,richardson2023}. These two developments raise a question at the heart of sustainable development. Is meeting basic human needs inherently tied to large environmental pressures, or can social progress be achieved within the limits of a stable planet?

\citet{raworth2012,raworth2017} gave the question a simple visual form, the ``doughnut''. Its inner ring is a social foundation, the minimum level of well-being to which everyone is entitled. Its outer ring is an ecological ceiling set by the planetary boundaries. Between the two lies a ``safe and just space'' in which human needs are met without overshooting the planet's limits. The framework does not say how to get there, but it makes the goal concrete. It also makes the goal measurable country by country, once global boundaries are translated into national shares and social thresholds are defined in observable terms.

This paper uses that idea to study how countries moved through the space between the social foundation and the climate boundary from \nFirstYear{} to \nLastYear. The panel covers \nNCountries{} countries. The main analyses use the \nNCons{} of them for which consumption-based emissions are available, and production-based emissions extend the key checks to all \nNCountries. It focuses on the climate boundary because climate change is one of the two ``core'' boundaries, whose transgression alone could shift the Earth system into a new state \citep{steffen2015}. Climate is also where the best long-run national data exist. The climate boundary is translated into an equal per-person share of a carbon budget consistent with the Paris Agreement's 2\textdegree C goal: \nBoundary~\tpc{} \citep{oneill2018}. Emissions are measured on a consumption basis, which assigns them to the country where goods and services are finally used \citep{peters2011,friedlingstein2026}. On the social side, four thresholds with near-complete data are used, taken from the literature on national doughnut accounting: life expectancy of at least 74 years, under-5 mortality of at most 25 per 1{,}000 live births, electricity access of at least 95\%, and employment of at least 94\% of the labour force \citep{fanning2022,un2015sdg}.

\subsection*{What is already known}

Two strands of research bear directly on this question. The first measures how far countries are from social thresholds and per-capita biophysical boundaries. \citet{oneill2018} assessed about 150 countries around 2011 on eleven social indicators and seven biophysical ones. They found that no country met the basic needs of its residents at a globally sustainable level of resource use. \citet{fanning2022} extended the analysis to a 1992--2015 panel of more than 140 countries and found that countries were crossing biophysical boundaries faster than they were achieving social thresholds. \citet{fanning2025} have since updated the doughnut to 2000--2022 with 35 indicators, reporting global trends and comparisons between groups of countries. The \citet{undp2020} now publishes a version of its Human Development Index adjusted for planetary pressures.

The second strand studies how human well-being relates to energy use and emissions. Life expectancy and other measures of well-being rise steeply with energy and carbon at low levels and then level off \citep{steinberger2010,lamb2015}. The carbon intensity of well-being, emissions per unit of life expectancy, has fallen in some world regions but not in others \citep{jorgenson2014}. Countries achieve similar outcomes with very different levels of energy use, depending on how basic services are provided \citep{vogel2021}. \citet{porcher2026} describe a ``champagne curve'' in which human development rises with consumption-based emissions at low levels and then flattens, so that beyond a point additional emissions add little to well-being. A related literature asks whether economic growth can be decoupled from emissions, and finds that absolute decoupling has occurred in some high-income countries, but far too slowly to meet climate targets \citep{haberl2020,vogel2023}.

\subsection*{What this paper adds}

This paper brings the two strands together at the level of individual countries, and contributes in three ways.

First, it introduces a simple and transparent summary of the relationship between emissions and social attainment across countries: the \emph{carbon requirement} of meeting the social foundation. This is the level of \COtwo{} per person at which, given the observed relationship across countries in a given year, half of countries meet the foundation. It is estimated for every year from \nShareYearStart{} to \nLastYear{} with country-cluster bootstrap confidence intervals, so that its change over time can be measured. The carbon requirement is descriptive. It summarises the cross-country distribution and is not an estimate of the emissions that any particular country needs to provide basic services. This distinction is emphasised throughout.

Second, it traces every country that started the period below the social foundation and later met it, and asks on which side of the climate boundary the crossing happened. This turns the static question ``are there countries in the safe and just space?'' into a dynamic one: does the historical path into the social foundation run through the safe and just space, or around it?

Third, it compares the rates at which countries actually decoupled emissions from income with the rates each would need to reach the climate boundary by 2050. \citet{vogel2023} made this comparison for 11 high-income countries. This paper extends it to all \nNCons{} countries with consumption-based emissions data, using a per-person boundary. It also reports scenarios in which countries cut emissions faster than they have so far, to show how large the gap is.

Methodologically, the analysis is designed to be fully reproducible. Every number, table and figure in the paper is generated by a single analysis pipeline from raw data files whose checksums are recorded. Numbers quoted in the text are inserted automatically from the pipeline's output rather than typed by hand. The pipeline also records diagnostics that are often left implicit, including failed statistical fits, the number of observations whose classification depends on interpolated values, and the characteristics of countries excluded for lack of data.

\subsection*{Main findings}

The analysis yields four main findings.
\begin{enumerate}[label=(\roman*)]
\item \textbf{The safe and just space has stayed empty.} None of the \nNCons{} countries with consumption-based emissions data met all four social thresholds while emitting within the climate boundary in any year between \nFirstYear{} and \nLastYear. With production-based emissions, only \nProdSafeCountriesAll{} did so in some years, and neither did once emissions embodied in imports are counted. All \nNCrossers{} countries that reached the social foundation during the period did so while overshooting the boundary, at a median of \nCrossCarbonMedian~t of \COtwo{} per person. The same holds for all \nNCrossersSust{} countries that met the foundation for at least three consecutive years.
\item \textbf{The carbon cost of meeting basic needs is falling, but remains far above the boundary.} The carbon requirement fell from about \nReqStart~t per person in \nShareYearStart{} to \nReqEnd~t in \nLastYear, an average decline of \nReqTrend\% a year. Yet in \nLastYear{} it was still \nReqEndMultiple{} times the boundary. For the two health thresholds alone it fell to \nReqHealthEnd~t, much closer to the boundary. The level of the requirement depends on how the social foundation is defined; its decline, and its position above the boundary, do not.
\item \textbf{Decoupling has spread but is far too slow.} The number of countries whose consumption-based emissions fell while their incomes grew rose from \nDecAbsEarly{} in 2001--2010 to \nDecAbsLate{} in 2011--2022. But only \nDecAboveOnTrack{} of the \nDecAbsAbove{} decouplers still above the boundary cut emissions fast enough to reach it by 2050. Even if all of them matched the fastest rate any of them achieved, only \nScenFastest{} would arrive in time.
\item \textbf{The configurations are persistent, and the typical path runs through overshoot.} Between 2001--2010 and 2011--2022, \nPersistTwo\% of countries stayed in the same configuration. Movement was almost entirely from social shortfall towards meeting the foundation, and from within the boundary to above it.
\end{enumerate}

These results confirm, with data to \nLastYear{} and consumption-based emissions, that no country in the sample has historically met these basic needs within its share of the climate boundary. They add evidence that the carbon cost of doing so has fallen substantially, and quantify how far current rates of decarbonisation fall short of what the boundary requires.

\subsection*{Structure of the paper}

\Cref{sec:literature} reviews the planetary boundaries, social foundations and decoupling literatures in more detail. \Cref{sec:framework} sets out the conceptual framework: the four configurations, the carbon requirement and the expectations that guide the analysis. \Cref{sec:data} describes the data, the construction of the sample and the treatment of missing values. \Cref{sec:methods} describes the methods, and \cref{sec:results} presents the results. \Cref{sec:discussion} discusses their interpretation, implications and limitations, and \cref{sec:conclusion} concludes. The appendices give the full country-level results and additional estimates.
